Задача упорядочивания данных является одной из фундаментальных в области
информатики и программирования. Алгоритмы сортировки применяются повсеместно"---
от баз данных и операционных систем до компиляторов и графических движков.
Эффективность выбранного алгоритма напрямую влияет на производительность
программной системы в целом.

Среди множества существующих алгоритмов сортировки особое место занимают
два: быстрая сортировка (\textit{Quicksort}), предложенная Ч.\,А.\,Р.~Хоаром
в 1962~году~\cite{hoare1962}, и сортировка слиянием (\textit{Merge sort}),
впервые описанная Джоном фон Нейманом в 1945~году. Оба алгоритма относятся
к классу <<разделяй и властвуй>> и широко используются на практике.

\textbf{Цель работы}"--- провести теоретический и практический анализ
алгоритмов быстрой сортировки и сортировки слиянием, выявить их сильные
и слабые стороны, а также определить условия предпочтительного применения
каждого из них.

Для достижения поставленной цели необходимо решить следующие задачи:
\begin{enumerate}
    \item изучить теоретические основы алгоритмов быстрой сортировки
          и сортировки слиянием;
    \item реализовать оба алгоритма на языках программирования C++ и Python;
    \item провести сравнительный анализ временной сложности и использования
          памяти;
    \item сформулировать рекомендации по выбору алгоритма.
\end{enumerate}
